\section{Results: the coordination tax}
\label{sec:results}

We measured the cost as a paired ensemble: 32 seeded galaxies of 300 stars each, every one
filled twice, once with perfect information and once with the light-speed gate of
(\ref{eq:gate}), so that the only difference between the two runs of a pair is the delay.
The primary quantity is the relative slowdown in the time to settle the entire field,
$t_{100}$, per seed; because $t_{100}$ is a tail statistic dominated by the last few stars,
we also report the penalty at earlier coverage fractions and across a range of field sizes.
Reporting the distribution across the ensemble rather than a single run follows the seeded
Monte Carlo practice of the field \cite{forgan-2009}.

\begin{table}[!t]
  \centering
  \caption{Fill-100\% penalty from light-speed lag over the 32-seed paired ensemble at
    300 stars: median, interquartile range, and a seeded percentile-bootstrap 95\%
    confidence interval on the median. The last column is a two-sided sign test over the
    paired seeds; the slingshot penalty is positive in all 32, while powered is null
    (positive in only 6 of the 11 seeds that differ at all).}
  \label{tab:penalty}
  \begin{tabular}{lcccc}
    \toprule
    Flight policy & Median (\%) & IQR (\%) & 95\% CI (\%) & $p$ \\
    \midrule
    Powered nearest-neighbour & 0.0  & [0.0, 0.0]   & [0.0, 0.0]   & 1.0 \\
    Nearest-star slingshot    & 30.3 & [20.0, 38.5] & [23.5, 37.5] & $<10^{-9}$ \\
    Maximum-boost slingshot   & 51.4 & [46.0, 54.0] & [47.2, 53.7] & $<10^{-9}$ \\
    \bottomrule
  \end{tabular}
\end{table}

\begin{table}[!t]
  \centering
  \caption{The mechanism, as medians over the 32 light-speed runs. Effective speed is the
    mean speed at which probes are launched; $\Lambda_{\text{eff}} = v_{\text{eff}}/c$. Hop
    length is the mean distance of a wasted (lost-race) trip and of a winning (settling)
    trip. The two slingshot policies differ in $\Lambda_{\text{eff}}$ by about
    $1.4\times$ but in wasted-hop length by about $1.85\times$, and the penalty in
    Table~\ref{tab:penalty} tracks the latter.}
  \label{tab:mech}
  \begin{tabular}{lcccc}
    \toprule
    Flight policy & $v_{\text{eff}}$ (km/s) & $\Lambda_{\text{eff}}$ & Wasted (pc) & Settled (pc) \\
    \midrule
    Powered nearest-neighbour & 9     & $3.0\times10^{-5}$ & 1.62 & 0.83 \\
    Nearest-star slingshot    & 2{,}734 & $9.1\times10^{-3}$ & 1.44 & 0.83 \\
    Maximum-boost slingshot   & 3{,}897 & $1.3\times10^{-2}$ & 2.67 & 2.36 \\
    \bottomrule
  \end{tabular}
\end{table}

The medians and their spread are given in Table~\ref{tab:penalty} and shown per seed in
Fig.~\ref{fig:slowdown}. Powered flight pays essentially nothing; nearest-star slingshot
loses about a third of its time; maximum-boost slingshot loses about half. The two slingshot
penalties are not marginal: each is positive in all 32 seeds, and the bootstrap confidence
intervals on the medians (Table~\ref{tab:penalty}) are well clear of zero and of each other.

The central finding is that the penalty is \emph{not} governed by $\Lambda = v/c$ from
(\ref{eq:lambda}). Table~\ref{tab:mech} makes the case directly. The two slingshot policies
reach comparable effective speeds, $2{,}734$ against $3{,}897$ km/s, so their
$\Lambda_{\text{eff}}$ differ by only about $1.4\times$; yet their penalties differ by about
$1.7\times$. If $\Lambda$ set the cost, the faster policy would pay in that same $1.4$-to-$1$
proportion, and the powered policy, though slowest, would still pay something. Neither holds.
What does track the penalty is hop non-locality. A probe flying to the \emph{nearest} star
that loses a race has lost almost nothing: the next unsettled star is right beside it, so the
wasted trip and the recovery are both short. A probe chasing the \emph{biggest boost}
deliberately reaches past its neighbours to a distant star, so when it loses that race both
the wasted trip and the re-targeting detour are long. The simulation confirms the mechanism
in the hop lengths themselves: the mean wasted trip is $2.67$ pc for maximum-boost against
$1.44$ pc for nearest-star, and its winning trips are nearly three times longer too
(Table~\ref{tab:mech}). So $\Lambda$ decides only \emph{whether} lag can bite: powered flight,
with $\Lambda \approx 3\times10^{-5}$, is immune not because its wasted hops are short (they
are not, at $1.62$ pc) but because it is far too slow for any race to matter. Given that a
policy is fast enough for lag to bite, the \emph{non-locality of its hops} sets how hard.
Fig.~\ref{fig:settlement} shows the effect on a single representative galaxy for the
nearest-star policy: the light-speed run reaches full settlement at about 110{,}000 years
against 80{,}000 for the perfect-information run on the same seed, a right-shift of roughly a
third visible as a lag in the settlement curve, not a change in its final height.

The penalty is not an artifact of the fragile $t_{100}$ tail. Fig.~\ref{fig:coverage} reports
the median penalty at coverage fractions from $t_{25}$ to $t_{100}$: for nearest-star flight
it rises steadily from about 14\% at quarter-coverage to about 30\% at full coverage, and for
maximum-boost it is already about 26\% at quarter-coverage and stays near or above 50\%
thereafter. The shift is present throughout the fill, and grows as the front reaches farther
from the origin, exactly as the locality mechanism predicts; it is not a single unlucky last
star. Powered flight stays at zero at every fraction.

Nor is it an artifact of the small $300$-star box, the concern that matters most for
extrapolating toward galactic saturation. Fig.~\ref{fig:finitesize} sweeps the system size
across a four-fold range: the powered penalty stays near zero, the nearest-star penalty stays
near a third, and the maximum-boost penalty holds near or slightly above a half, firming up
rather than washing out as the field grows and hops lengthen. The maximum-boost policy scans
its nearest neighbours at every launch and so costs $O(N^2)$, which bounds the range we sweep;
within that range the penalties are stable, and the trend is toward more penalty at larger
scale, not less.

Finally, what does \emph{not} happen: every case still fills the field to full coverage. This
is a property of the model, not a result to be discovered. Re-targeting guarantees that a
probe which loses a race simply picks another star, and with no loss term the reachable field
must eventually close; light-speed lag makes the swarm slower and more wasteful but leaves no
permanent unsettled remainder. It is worth stating precisely because it marks the boundary of
what delay alone can do: a permanent settled fraction below one, an ``Aurora'' steady state in
which some region stays forever empty \cite{carroll-nellenback-2019}, would require settled
sites to \emph{die}, and delay supplies no such term. Lag is a tax on speed, not a barrier to
completion.

\begin{figure}[!t]
  \centering
  \includegraphics[width=\columnwidth]{fig_slowdown_by_policy.pdf}
  \caption{Distribution of the per-seed fill-100\% penalty imposed by light-speed-limited
    coordination, over the fixed 32-seed ensemble at 300 stars, for the three flight policies
    (box: interquartile range and median; points: all 32 individual seeds). Powered
    nearest-neighbour flight pays essentially no penalty, because a lost race is recovered at
    the star next door, while the long-range maximum-boost policy fills the field materially
    later, because wasted trips launched from stale views become long detours.}
  \label{fig:slowdown}
\end{figure}

\begin{figure}[!t]
  \centering
  \includegraphics[width=\columnwidth]{fig_penalty_by_coverage.pdf}
  \caption{Median fill-time penalty at coverage fractions from $t_{25}$ to $t_{100}$, over the
    32-seed ensemble, for the three policies. The slowdown is present across the whole fill,
    not only in the $t_{100}$ tail, and grows as the front reaches farther from the origin;
    powered flight stays at zero throughout.}
  \label{fig:coverage}
\end{figure}

\begin{figure}[!t]
  \centering
  \includegraphics[width=\columnwidth]{fig_finite_size.pdf}
  \caption{Median fill-100\% penalty versus system size for the three policies (line: median;
    shaded band: interquartile range over 8 seeds per point). Across a four-fold increase in
    the number of stars the penalties are stable, holding near 0\%, a third, and a half
    respectively, and firm up rather than wash out. The $O(N^2)$ cost of the maximum-boost
    scan bounds the range swept.}
  \label{fig:finitesize}
\end{figure}

\begin{figure}[!t]
  \centering
  \includegraphics[width=\columnwidth]{fig_settlement_curves.pdf}
  \caption{Fraction of the reachable field settled versus year for the nearest-star slingshot
    policy on a single representative galaxy, under perfect global information (solid) and
    under light-speed-limited coordination (dashed). The light-speed curve is shifted later,
    reaching full settlement at about 110{,}000 years versus 80{,}000, a delay attributable to
    information lag alone; both curves reach the same final height.}
  \label{fig:settlement}
\end{figure}
